\documentclass[a4paper]{book}

\setlength{\headheight}{14.5pt}

% --- 1. ENCODAGE ET LANGUE (À mettre en premier) ---
\usepackage[T1]{fontenc}
\usepackage[utf8]{inputenc}
\usepackage[french]{babel}

% --- 2. POLICES ET SYMBOLES ---
% fourier : police élégante pour les maths (lmodern supprimé car doublon)
\usepackage{fourier}
\usepackage{amsfonts, amssymb, amsthm}
\usepackage[right]{eurosym}
\usepackage[np]{numprint}
\usepackage[nointegrals]{wasysym}

% --- 3. GÉOMÉTRIE ET MISE EN PAGE ---
\usepackage[left=1cm,right=1cm,top=1.5cm,bottom=1.5cm]{geometry}
\usepackage[dvipsnames,svgnames]{xcolor}
\usepackage[most,skins,breakable]{tcolorbox}
\usepackage{tagging}

% --- 4. GRAPHISMES ET DIAGRAMMES ---
\usepackage{graphicx}
\usepackage{pgfplots}
\pgfplotsset{compat=1.18}
\usetikzlibrary{babel, shapes, fit} % 'babel' évite les bugs avec TikZ

% --- 5. TABLEAUX ---
\usepackage{array, booktabs, multirow, tabularx, longtable}

% --- 6. STRUCTURE ET TITRES ---
% titlesec est plus moderne que sectsty, on garde uniquement celui-ci
\usepackage{titlesec}
\usepackage{enumitem}
\usepackage[normalem]{ulem} % 'normalem' évite que \emph devienne souligné

% --- 7. HYPERLIENS (toujours en dernier ou presque) ---
\usepackage{hyperref}

\setlength{\parindent}{0pt}
\frenchbsetup{StandardItemLabels=true}

% Le package ultime
\usepackage{ProfCollege}


% =================================================================
% COMMANDES GÉNÉRALES
% =================================================================

% Une ligne de points (usage : \points[15])
\newcommand{\points}[1][1]{%
    \foreach \n in {1,...,#1} { \ldots}%
}

% Plusieurs lignes de points (usage : \lignesPoints{3})
\newcommand{\lignesPoints}[1]{%
    \foreach \n in {1,...,#1} {%
        \dotfill \newline
    }%
}


% =================================================================
% ZONES RÉPONSE
% =================================================================

% Zone réponse universelle (usage : \begin{ZoneReponse}{4cm})
\newtcolorbox{ZoneReponse}[1]{
    enhanced,
    colback=white,
    colframe=gray!50,
    arc=2mm,
    boxrule=0.5mm,
    width=\textwidth,
    height=#1,
    valign=top,
    fonttitle=\bfseries,
    colbacktitle=gray!20,
    coltitle=black,
    attach title to upper,
    after skip=10pt plus 2pt minus 2pt
}


% =================================================================
% COMMANDES MATHÉMATIQUES COLLÈGE
% =================================================================

% Ensembles de nombres
\newcommand{\N}{\ensuremath{\mathbf{N}}}
\newcommand{\Z}{\ensuremath{\mathbf{Z}}}
\newcommand{\D}{\ensuremath{\mathbf{D}}}
\newcommand{\Q}{\ensuremath{\mathbf{Q}}}
\newcommand{\R}{\ensuremath{\mathbf{R}}}

% Valeur absolue
\newcommand{\abs}[1]{\left\lvert #1 \right\rvert}

% Arc géométrique (usage : $\arc{AB}$)
% Note : utilise la fonte tipa pour le symbole arc
\makeatletter
\DeclareFontFamily{U}{tipa}{}
\DeclareFontShape{U}{tipa}{m}{n}{<->tipa10}{}
\newcommand{\arc@char}{{\usefont{U}{tipa}{m}{n}\symbol{62}}}
\newcommand{\arc}[1]{\mathpalette\arc@arc{#1}}
\newcommand{\arc@arc}[2]{%
  \sbox0{$\m@th#1#2$}%
  \vbox{\hbox{\resizebox{\wd0}{\height}{\arc@char}}\nointerlineskip\box0}%
}
\makeatother

% Comparaison : boîte vide entre deux nombres pour écrire le symbole
% Usage : \compare{102}{120}  ou  \compare[0.5]{\np{132999999}}{\np{133000000}}
% Argument optionnel : ignoré (compatibilité ascendante), la boîte s'adapte
\newcommand{\compare}[3][0]{%
    #2 \kern0.4em
    \tikz[baseline=-0.6ex]
        \node[rectangle, draw=gray!60, thick, dotted, minimum height=0.6cm,
              minimum width=0.6cm, inner sep=2pt] {};%
    \kern0.4em #3%
}

% Encadrement : deux boîtes vides autour d'un nombre
% Usage : \encadre{\np{1389}}  ou  \encadre[1.8]{\np{652176}}
% Argument optionnel : largeur des boîtes en cm (défaut 1.5)
\newcommand{\encadre}[2][1.5]{%
    \tikz[baseline=-0.6ex]
        \node[rectangle, draw=gray!60, thick, dotted, minimum height=0.6cm,
              minimum width=#1cm, inner sep=2pt] {};%
    \kern0.3em $<$ #2 $<$ \kern0.3em
    \tikz[baseline=-0.6ex]
        \node[rectangle, draw=gray!60, thick, dotted, minimum height=0.6cm,
              minimum width=#1cm, inner sep=2pt] {};%
}

% Diagramme de proportionnalité simple
% Usage : \diagrammesimple{A}{B}{×3}
\newcommand{\diagrammesimple}[3]{%
\begin{tikzpicture}[baseline=(A.base)]
    \node(A) at (0,0) {#1};
    \node(B) at (5,0) {#2};
    \draw[->, dashed, line width=0.8mm, green!50!black]
        (A) to[bend left=20] node[midway, above=2pt] {#3} (B);
\end{tikzpicture}}

% Diagramme avec flèche aller-retour
% Usage : \diagrammesimplereciproque{A}{B}{×3}{÷3}
\newcommand{\diagrammesimplereciproque}[4]{%
\begin{tikzpicture}[baseline=(A.base)]
    \node(A) at (0,0) {#1};
    \node(B) at (5,0) {#2};
    \draw[->, dashed, line width=0.8mm, green!50!black]
        (A) to[bend left=20] node[midway, above=2pt] {#3} (B);
    \draw[->, dashed, line width=0.8mm, blue!50!black]
        (B) to[bend left=20] node[midway, below=2pt] {#4} (A);
\end{tikzpicture}}
\newcommand{\maClasse}{3e}

%% --- PACKAGES PROPRES AUX COURS ---
\usepackage{fancyhdr}
\usepackage{lastpage}
\usepackage[column=X]{cellspace}

%% --- VARIABLES ---
\newcommand{\monTitre}{Titre du cours} % À redéfinir dans chaque fichier : \renewcommand{\monTitre}{Mon titre}

%% --- FORMAT DES TITRES (titlesec) ---

\titleformat{\chapter}[frame]
{\color{red!85!white}\normalsize}
{\filright\sffamily\Large\enspace Chapitre \thechapter\enspace}
{8pt}
{\color{red!85!white}\rule{0pt}{30pt}\sffamily\Huge\bfseries\filcenter}
\titlespacing{\chapter}{0pc}{0pc}{2pc}

\titleformat{\section}
{\color{green!55!black}\normalfont\Large\bfseries}
{\thesection}
{1em}
{}

% 3. Subsections (Cyan)
% J'ai retiré le \hspace{1em} qui poussait le texte
\titleformat{\subsection}[block]
{\color{cyan!65!black}\normalfont\large\bfseries}
{\thesubsection}
{1em}
{}
% J'ai mis le premier paramètre à 0pt pour supprimer l'indentation de bloc
\titlespacing{\subsection}{0pt}{1.5ex}{1ex}

% 4. Subsubsections (Violet)
\titleformat{\subsubsection}[block]
{\color{violet}\normalfont\normalsize\bfseries}
{\thesubsubsection}
{1em}
{}
\titlespacing{\subsubsection}{0pt}{1ex}{1ex}

%% --- NUMÉROTATION ---
\renewcommand{\thesection}{\Roman{section}.}
\renewcommand{\thesubsection}{\arabic{subsection})}
\renewcommand{\thesubsubsection}{\alph{subsubsection})}
\setcounter{secnumdepth}{3}

%% --- EN-TÊTES ET PIEDS DE PAGE ---
\pagestyle{fancy}
\renewcommand\headrulewidth{1pt}
\fancyhead[LO,RE]{\maClasse}
\fancyhead[RO,LE]{\monTitre}
\renewcommand\footrulewidth{1pt}
\fancyfoot[LO,RE]{Année 2025-2026}
\fancyfoot[RO,LE]{\textbf{Page \thepage/\pageref{LastPage}}}

%% --- RÉGLAGES PARAGRAPHES ---
\setlength{\cellspacebottomlimit}{4pt}
\setlength{\cellspacetoplimit}{4pt}
\setlength{\parindent}{2pt}
\setlength{\tabcolsep}{0.5cm}

%% --- COMPTEURS ET BOITES TCOLORBOX ---

%% --- COMPTEURS AVEC RÉINITIALISATION À CHAQUE CHAPITRE ---

% On ajoute [chapter] à la fin pour que le compteur reparte à 0 à chaque \chapter
\newcounter{dfn}[chapter]
\newcounter{thm}[chapter]
\newcounter{cvn}[chapter]
\newcounter{prp}[chapter]
\newcounter{ex}[chapter]
\newcounter{mtd}[chapter]
\newcounter{rq}[chapter]
\newcounter{demo}[chapter]
\newcounter{appli}[chapter]
\newcounter{exo}[chapter]
\newcounter{lem}[chapter]

%% --- BLOC COMPLET DES BOITES TCOLORBOX ---

% --- DÉFINITIONS (Bleu) ---
\newtcolorbox{definition}{breakable,colback=white!6!white,colframe=blue!85!white,fonttitle=\bfseries,title=Définition \thedfn, code={\stepcounter{dfn}}, arc=3mm}
\newtcolorbox{definitionTitre}[1]{breakable,colback=white!6!white,colframe=blue!85!white,fonttitle=\bfseries,title=Définition \thedfn\ - #1, code={\stepcounter{dfn}}, arc=3mm}
\newtcolorbox{definitions}{breakable,colback=white!6!white,colframe=blue!85!white,fonttitle=\bfseries,title=Définitions \thedfn, code={\stepcounter{dfn}}, arc=3mm}

% --- PROPRIÉTÉS & THÉORÈMES (Rouge/Orange) ---
\newtcolorbox{theoreme}{breakable,colback=red!6!white,colframe=red!85!white,fonttitle=\bfseries,title=Théorème \thethm, code={\stepcounter{thm}}, arc=3mm}
\newtcolorbox{theoremeTitre}[1]{breakable,colback=red!6!white,colframe=red!85!white,fonttitle=\bfseries,title=Théorème \thethm\ - #1, code={\stepcounter{thm}}, arc=3mm}
\newtcolorbox{propriete}{breakable,colback=red!6!white,colframe=red!85!white,fonttitle=\bfseries,title=Propriété \theprp, code={\stepcounter{prp}}, arc=3mm}
\newtcolorbox{proprieteTitre}[1]{breakable,colback=red!6!white,colframe=red!85!white,fonttitle=\bfseries,title=Propriété \theprp\ - #1, code={\stepcounter{prp}}, arc=3mm}
\newtcolorbox{proprietes}{breakable,colback=red!6!white,colframe=red!85!white,fonttitle=\bfseries,title=Propriétés \theprp, code={\stepcounter{prp}}, arc=3mm}
\newtcolorbox{lemme}{breakable,colback=red!6!white,colframe=orange!85!white,fonttitle=\bfseries,title=Lemme \theprp, code={\stepcounter{prp}}, arc=3mm}
\newtcolorbox{convention}{breakable,colback=orange!6!white,colframe=orange!85!white,fonttitle=\bfseries,title=Convention \thecvn, code={\stepcounter{cvn}}, arc=3mm}

% --- EXEMPLES (Vert) — argument optionnel pour le titre ---
% Usage : \begin{exemple}             -> "Exemple 1"
%         \begin{exemple}[Mon titre]  -> "Exemple 1 - Mon titre"
\newtcolorbox{exemple}[1][]{
  enhanced,
  breakable,
  nobeforeafter,
  % --- ON FORCE TOUT À ZÉRO ---
  boxrule=0pt,        % Épaisseur du cadre global à 0
  titlerule=0pt,      % Supprime la ligne sous le titre
  sharp corners,      % Important pour les coupures propres
  % ----------------------------
  colback=white,
  colframe=white,     % On met le cadre en blanc (invisible) au cas où
  frame empty,        % On cache le cadre
  %
  borderline west={0.5mm}{0mm}{green!55!black}, % Seule la barre verte reste
  %
  % Gestion de la coupure
  toprule at break=0pt,
  bottomrule at break=0pt,
  %
  title={Exemple \theex\ifx&#1&\else\ - #1\fi},
  code={\stepcounter{ex}},
  attach boxed title to top left = {xshift = -3mm},
  boxed title style={size=small,colback=green!55!black,colframe=green!55!black, arc=3mm}
}

% --- REMARQUES (Orange) ---
\newtcolorbox{remarque}{enhanced, breakable, frame empty, nobeforeafter, borderline west={0.5mm}{0mm}{orange!85!white}, title=Remarque \therq, colback=white, code={\stepcounter{rq}}, attach boxed title to top left = {xshift = -3mm}, boxed title style={size=small,colback=orange!85!white,colframe=orange!85!white, arc=3mm}}
\newtcolorbox{remarques}{enhanced, breakable, frame empty, nobeforeafter, borderline west={0.5mm}{0mm}{orange!85!white}, title=Remarques \therq, colback=white, code={\stepcounter{rq}}, attach boxed title to top left = {xshift = -3mm}, boxed title style={size=small,colback=orange!85!white,colframe=orange!85!white, arc=3mm}}
\newtcolorbox{remarqueTitre}[1]{enhanced, breakable, frame empty, nobeforeafter, borderline west={0.5mm}{0mm}{orange!85!white}, title=Remarque \therq\ - #1, colback=white, code={\stepcounter{rq}}, attach boxed title to top left = {xshift = -3mm}, boxed title style={size=small,colback=orange!85!white,colframe=orange!85!white, arc=3mm}}

% --- MÉTHODES & APPLICATIONS (Bleu foncé) ---
\newtcolorbox{methode}{enhanced, breakable, frame empty, nobeforeafter, borderline west={0.5mm}{0mm}{blue!75!black}, title=Méthode \themtd, colback=white, code={\stepcounter{mtd}}, attach boxed title to top left = {xshift = -3mm}, boxed title style={size=small,colback=blue!75!black,colframe=blue!75!black, arc=3mm}}
\newtcolorbox{methodes}{enhanced, breakable, frame empty, nobeforeafter, borderline west={0.5mm}{0mm}{blue!75!black}, title=Méthodes \themtd, colback=white, code={\stepcounter{mtd}}, attach boxed title to top left = {xshift = -3mm}, boxed title style={size=small,colback=blue!75!black,colframe=blue!75!black, arc=3mm}}
\newtcolorbox{methodeTitre}[1]{enhanced, breakable, frame empty, nobeforeafter, borderline west={0.5mm}{0mm}{blue!75!black}, title=Méthode \themtd\ - #1, colback=white, code={\stepcounter{mtd}}, attach boxed title to top left = {xshift = -3mm}, boxed title style={size=small,colback=blue!75!black,colframe=blue!75!black, arc=3mm}}
\newtcolorbox{application}{enhanced, breakable, frame empty, nobeforeafter, borderline west={0.5mm}{0mm}{orange!75!black}, title=Application \themtd, colback=white, code={\stepcounter{mtd}}, attach boxed title to top left = {xshift = -3mm}, boxed title style={size=small,colback=orange!75!black,colframe=orange!75!black, arc=3mm}}

% --- DIVERS (Exercices, Python, Conséquences) ---
\newtcolorbox{exercice}{enhanced, breakable, frame empty, nobeforeafter, borderline west={0.5mm}{0mm}{green!55!black}, title=Exercice n°\theexo, colback=white, code={\stepcounter{exo}}, attach boxed title to top left = {xshift = -3mm}, boxed title style={size=small,colback=green!55!black,colframe=green!55!black, arc=3mm}}
\newtcolorbox{exerciceTitre}[1]{enhanced, breakable, frame empty, borderline west={0.5mm}{0mm}{green!55!black}, title=Exercice n°\theexo\ - #1, colback=white, code={\stepcounter{exo}}, attach boxed title to top left = {xshift = -3mm}, boxed title style={size=small,colback=green!55!black,colframe=green!55!black, arc=3mm}}
\newtcolorbox{consequences}{enhanced, breakable, frame empty, nobeforeafter, borderline west={0.5mm}{0mm}{ProcessBlue}, title=Conséquences, colback=white, attach boxed title to top left = {xshift = -3mm}, boxed title style={size=small,colback=ProcessBlue,colframe=ProcessBlue, arc=3mm}}
\newtcolorbox{programme}{enhanced, breakable, frame empty, nobeforeafter, borderline west={0.5mm}{0mm}{gray!55!black}, title=Programme Python, colback=white, attach boxed title to top left = {xshift = -3mm}, boxed title style={size=small,colback=gray!55!black,colframe=gray!55!black, arc=3mm}}


% Commande pour les démonstrations (incrémentation automatique)
\newcommand{\demonstration}{%
  \stepcounter{demo}%
  \tcbox[enhanced, 
         on line, 
         arc=2mm, 
         colback=gray!10!white, 
         colframe=gray!50!black, 
         size=small, 
         fontupper=\bfseries]
         {Démonstration \thedemo}%
  \enspace % Ajoute un petit espace après la boîte
}

%% --- FIN DU BLOC DES BOITES ---

% Activité (Numéro, Titre, Chapitre)
\newcommand{\enteteActivite}[3]{
    \noindent
    \begin{tabularx}{\textwidth}{@{}p{2cm} >{\centering\arraybackslash}X p{2cm}@{}}
        \textbf{Activité n°#1} & \textbf{#2} & \textbf{Chapitre #3} \\
    \end{tabularx}
}


% --- Commande pour les blocs mémo personnalisés ---
% Argument 1 : Couleur du cadre
% Argument 2 : Titre du bloc
% Argument 3 : Texte/Contenu
\newcommand{\blocMemo}[3]{%
    \begin{tcolorbox}[
        enhanced,
        breakable,
        title=#2,
        colback=#1!5,          % Fond très clair
        colframe=#1!75!black,  % Bordure
        fonttitle=\bfseries\sffamily,
        colbacktitle=#1!75!black,
        attach boxed title to top left={xshift=3mm, yshift=-3mm},
        boxed title style={size=small, arc=2mm, colframe=#1!75!black},
        arc=3mm,               % Bel arrondi
        boxrule=0.6mm,         % Trait un peu plus épais
        after skip=15pt,
        before skip=10pt
    ]
        \vspace{2mm} % Petit espace pour ne pas coller au titre boxé
        #3
    \end{tcolorbox}
}


% --- Nouvel environnement Culture Maths ---
\newtcolorbox{cultureMaths}{
    enhanced, 
    breakable, 
    frame empty, 
    nobeforeafter, 
    borderline west={0.5mm}{0mm}{gray!60!black}, % Bordure grise comme ton "programme"
    title=Le saviez-vous ?, 
    colback=white, 
    attach boxed title to top left = {xshift = -3mm}, 
    boxed title style={size=small,colback=gray!60!black,colframe=gray!60!black, arc=3mm},
    coltitle=white,
    fonttitle=\bfseries\sffamily,
    after skip=15pt,
    before skip=10pt
}

\begin{document}

\setcounter{chapter}{1}
\renewcommand{\monTitre}{Les puissances}
\chapter{\monTitre}
\thispagestyle{fancy}


% ============================================================
\section{Exposant entier positif}
% ============================================================

\subsection{Définition et notation}

    \begin{definition}
        Pour tout nombre $a$ et tout entier $n \geq 2$, on appelle \textbf{puissance $n$-ième de $a$} le produit de $n$ facteurs égaux à $a$ :
        \[
            a^n = \underbrace{a \times a \times \cdots \times a}_{n \text{ facteurs}}
        \]
        Par convention : $a^1 = a$ \quad et \quad $a^0 = 1$ \quad (pour $a \neq 0$).
        
        \medskip
        $a$ s'appelle la \textbf{base} et $n$ s'\textbf{appelle l'exposant}.
    \end{definition}

    \begin{exemple}
        Calculer les puissances suivantes :
        \begin{enumerate}
            \item $3^4 = $
            \begin{ZoneReponse}{1cm}
            \end{ZoneReponse}

            \item $(-2)^3 = $
            \begin{ZoneReponse}{1cm}
            \end{ZoneReponse}

            \item $(-5)^2 = $
            \begin{ZoneReponse}{1cm}
            \end{ZoneReponse}

            \item $7^0 = $
            \begin{ZoneReponse}{1cm}
            \end{ZoneReponse}
        \end{enumerate}
    \end{exemple}

    \begin{remarque}
        \textbf{Attention au signe !} Il ne faut pas confondre $(-a)^n$ et $-a^n$ :
        \begin{itemize}
            \item $(-3)^2 = (-3) \times (-3) = 9$ \quad (on élève $-3$ au carré)
            \item $-3^2 = -(3 \times 3) = -9$ \quad (on prend l'opposé de $3^2$)
        \end{itemize}
        De plus, une puissance d'exposant \textbf{pair} d'un nombre négatif est \textbf{toujours positive}.
        Une puissance d'exposant \textbf{impair} d'un nombre négatif est \textbf{toujours négative}.
    \end{remarque}


\subsection{Propriétés de calcul}

    \begin{propriete}[Produit de puissances de même base]
        Pour tout nombre $a \neq 0$ et tous entiers $m$ et $n$ :
        \[
            a^m \times a^n = a^{m+n}
        \]
    \end{propriete}

    \begin{exemple}
        Simplifier en utilisant la propriété :
        \begin{enumerate}
            \item $5^3 \times 5^4 = $
            \begin{ZoneReponse}{1cm}
            \end{ZoneReponse}

            \item $(-2)^5 \times (-2)^2 = $
            \begin{ZoneReponse}{1cm}
            \end{ZoneReponse}
        \end{enumerate}
    \end{exemple}

    \begin{propriete}[Quotient de puissances de même base]
        Pour tout nombre $a \neq 0$ et tous entiers $m$ et $n$ :
        \[
            \frac{a^m}{a^n} = a^{m-n}
        \]
    \end{propriete}

    \begin{exemple}
        Simplifier en utilisant la propriété :
        \begin{enumerate}
            \item $\dfrac{7^6}{7^2} = $
            \begin{ZoneReponse}{1cm}
            \end{ZoneReponse}

            \item $\dfrac{3^5}{3^5} = $
            \begin{ZoneReponse}{1cm}
            \end{ZoneReponse}
        \end{enumerate}
    \end{exemple}

    \begin{propriete}[Puissance d'une puissance]
        Pour tout nombre $a$ et tous entiers $m$ et $n$ :
        \[
            \left(a^m\right)^n = a^{m \times n}
        \]
    \end{propriete}

    \begin{exemple}
        Simplifier :
        \begin{enumerate}
            \item $\left(4^3\right)^2 = $
            \begin{ZoneReponse}{1cm}
            \end{ZoneReponse}

            \item $\left(10^2\right)^4 = $
            \begin{ZoneReponse}{1cm}
            \end{ZoneReponse}
        \end{enumerate}
    \end{exemple}

    \begin{propriete}[Puissance d'un produit ou d'un quotient]
        Pour tous nombres $a$ et $b$ et tout entier $n$ :
        \[
            (a \times b)^n = a^n \times b^n \qquad \text{et} \qquad \left(\frac{a}{b}\right)^n = \frac{a^n}{b^n} \quad (b \neq 0)
        \]
    \end{propriete}

    \begin{exemple}
        Développer puis calculer :
        \begin{enumerate}
            \item $(3 \times 5)^2 = $
            \begin{ZoneReponse}{1cm}
            \end{ZoneReponse}

            \item $\left(\dfrac{2}{3}\right)^3 = $
            \begin{ZoneReponse}{1cm}
            \end{ZoneReponse}
        \end{enumerate}
    \end{exemple}


% ============================================================
\section{Exposant entier négatif}
% ============================================================

\subsection{Définition}

    \begin{definition}
        Pour tout nombre $a \neq 0$ et tout entier $n > 0$, $a^{-n}$ est l'inverse de  $a^n$, ce qui veut dire que:
        \[
            a^{-n} = \frac{1}{a^n}
        \]
        En particulier : $a^{-1} = \dfrac{1}{a}$.
    \end{definition}

    \begin{exemple}
        Écrire chaque puissance sous forme de fraction puis calculer :
        \begin{enumerate}
            \item $2^{-3} = $
            \begin{ZoneReponse}{1.5cm}
            \end{ZoneReponse}

            \item $10^{-2} = $
            \begin{ZoneReponse}{1.5cm}
            \end{ZoneReponse}

            \item $5^{-1} = $
            \begin{ZoneReponse}{1.5cm}
            \end{ZoneReponse}
        \end{enumerate}
    \end{exemple}

    \begin{remarque}
        Les puissances de $10$ négatives correspondent à des fractions décimales :
        \[
            10^{-1} = 0{,}1 \qquad 10^{-2} = 0{,}01 \qquad 10^{-3} = 0{,}001 \qquad \cdots
        \]
    \end{remarque}


\newpage


\subsection{Les propriétés restent valables}

    \begin{propriete}
        Toutes les propriétés vues en section~1 (produit, quotient, puissance d'une puissance…) restent valables pour des exposants entiers \textbf{relatifs} (positifs, nuls ou négatifs).
    \end{propriete}

    \begin{exemple}
        Simplifier et écrire sous la forme $a^n$ avec $n$ entier relatif :
        \begin{enumerate}
            \item $10^3 \times 10^{-5} = $
            \begin{ZoneReponse}{1.5cm}
            \end{ZoneReponse}

            \item $\dfrac{2^{-4}}{2^{-1}} = $
            \begin{ZoneReponse}{1.5cm}
            \end{ZoneReponse}

            \item $\left(3^{-2}\right)^3 = $
            \begin{ZoneReponse}{1.5cm}
            \end{ZoneReponse}
        \end{enumerate}
    \end{exemple}


% ============================================================
\section{Écriture scientifique}
% ============================================================

\subsection{Définition}

    \begin{definition}
        Un nombre est écrit en \textbf{notation scientifique} (ou écriture scientifique) lorsqu'il est de la forme :
        \[
            a \times 10^n
        \]
        où $n$ est un entier relatif et $a$ est un nombre décimal tel que $1 \leq a < 10$.
    \end{definition}

    \begin{exemple}
        Parmi les écritures suivantes, indiquer celles qui sont des notations scientifiques et corriger les autres :
        \begin{enumerate}
            \item $3{,}5 \times 10^4$
            \begin{ZoneReponse}{0.8cm}
            \end{ZoneReponse}

            \item $15 \times 10^3$
            \begin{ZoneReponse}{0.8cm}
            \end{ZoneReponse}

            \item $0{,}8 \times 10^{-2}$
            \begin{ZoneReponse}{0.8cm}
            \end{ZoneReponse}

            \item $7 \times 10^{-5}$
            \begin{ZoneReponse}{0.8cm}
            \end{ZoneReponse}
        \end{enumerate}
    \end{exemple}


\subsection{Convertir en notation scientifique}

    \begin{methode}
        Pour écrire un nombre en notation scientifique :
        \begin{enumerate}
            \item On repère le premier chiffre significatif (le premier chiffre le plus à gauche différent de $0$).
            \item On place la virgule juste après ce chiffre pour obtenir $a$ avec $1 \leq a < 10$.
            \item On détermine $n$ : c'est le nombre de décimales dont on a \textbf{déplacé} la virgule.
                \begin{itemize}
                    \item Si on déplace vers la \textbf{gauche} (nombre $\geq 10$) : $n$ \textbf{est strictement positif}.
                    \item Si on déplace vers la \textbf{droite} (nombre $< 1$) : $n$ \textbf{est strictement négatif}.
                \end{itemize}
        \end{enumerate}
    \end{methode}

    \begin{exemple}
        Écrire les nombres suivants en notation scientifique :
        \begin{enumerate}
            \item $\np{384000}$
            \begin{ZoneReponse}{2cm}
            \end{ZoneReponse}

            \item $0{,}00047$
            \begin{ZoneReponse}{2cm}
            \end{ZoneReponse}

            \item $\np{5100000000}$ \quad (distance Terre-Jupiter en km, environ)
            \begin{ZoneReponse}{2cm}
            \end{ZoneReponse}

            \item $0{,}000\,000\,001$ \quad (taille d'une bactérie en m, environ)
            \begin{ZoneReponse}{2cm}
            \end{ZoneReponse}
        \end{enumerate}
    \end{exemple}


\subsection{Convertir depuis la notation scientifique}

    \begin{exemple}
        Écrire les nombres suivants sous forme décimale :
        \begin{enumerate}
            \item $6{,}02 \times 10^3$
            \begin{ZoneReponse}{1.5cm}
            \end{ZoneReponse}

            \item $3{,}1 \times 10^{-4}$
            \begin{ZoneReponse}{1.5cm}
            \end{ZoneReponse}
        \end{enumerate}
    \end{exemple}


\subsection{Opérations en notation scientifique}
 
    \begin{methode}
        Pour \textbf{multiplier} deux nombres en notation scientifique $(a \times 10^m) \times (b \times 10^n)$ :
        \begin{enumerate}
            \item On multiplie les parties décimales : $a \times b$.
            \item On additionne les exposants : $10^{m+n}$.
            \item Si le résultat n'est pas en notation scientifique, on ajuste.
        \end{enumerate}
        Pour \textbf{diviser} deux nombres en notation scientifique $\dfrac{a \times 10^m}{b \times 10^n}$ :
        \begin{enumerate}
            \item On divise les parties décimales : $a \div b$.
            \item On soustrait les exposants : $10^{m-n}$.
            \item Si le résultat n'est pas en notation scientifique, on ajuste.
        \end{enumerate}
    \end{methode}
 
    \begin{exemple}[Type Brevet – Calculs en notation scientifique]
        \begin{enumerate}
            \item Calculer $(137 \times 10^4) \times (2{,}5 \times 10^3)$ et donner le résultat en notation scientifique.
            \begin{ZoneReponse}{2.5cm}
            \end{ZoneReponse}

             \item Calculer $\dfrac{81 \times 10^7}{6 \times 10^{-2}}$ et donner le résultat en notation scientifique.
            \begin{ZoneReponse}{2.5cm}
            \end{ZoneReponse}
          
            \item La lumière parcourt environ $3 \times 10^8$ m par seconde. La distance de la Terre au Soleil est d'environ $1{,}5 \times 10^{11}$ m.
            Calculer le temps que met la lumière pour parcourir cette distance. Exprimer le résultat en secondes en notation scientifique, puis en minutes.
            \begin{ZoneReponse}{4cm}
            \end{ZoneReponse}
        \end{enumerate}
    \end{exemple}
 

\end{document}